\subsection{Multiplicity on a Graph} 
Consider the graph below of
\[
f(x)=\tfrac{1}{256}\left(x-4\right)^2\left(x+2\right)^3\left(x+6\right)
\]
\begin{center}\begin{tikzpicture}[x=0.5\linespace,y=0.5\linespace]
	\setgraphlarge
	\GraphSmall
	
	\draw[graph,domain=-3.13923:2.40972,samples=100] plot (2*\x,{0.25*(\x-2)^2*(\x+1)^3*(\x+3)});

\end{tikzpicture}\end{center}
\begin{itemize}
\item The zero \makeblanksmall has multiplicity \makeblanksmall. Near \blankpair, the graph looks like the graph of \makeblank[1in] (it's almost a \makeblank[1.5in])
\item The zero \makeblanksmall has multiplicity \makeblanksmall. Near \blankpair, the graph looks like the graph of \makeblank[1in]
\item The zero \makeblanksmall has multiplicity \makeblanksmall. Near \blankpair, the graph looks like the graph of \makeblank[1in] (it's almost a \makeblank[1.5in])
\end{itemize}
If a root $r$ has multiplicity $m$, then near that point, the graph looks like \makeblank[1in]